\section{Простые механизмы}
%%%%%%%%%%%%%%%%%%%%%%%%%%%%%%%%%%%%%%%%%%%%%%%%%%%%
\subsection{Рычаг}
%%%%%%%%%%%%%%%%%%%%%%%%%%%%%%%%%%%%%%%%%%%%%%%%%%%%
%%%%%%%%%%%%%%%%%%%%%%%%%%%%%%%%%%%%%%%%%%%%%%%%%%%%
%%%%%%%%%%%%%%%%%%%%%%%%%%%%%%%%%%%%%%%%%%%%%%%%%%%%
\z[{\begin{wrapfigure}[3]{r}{0.2\textwidth}%
		\vspace{-15pt}%
		\begin{flushright}%
 			\includegraphics[width=0.2\textwidth]{5.1.1}%
		\end{flushright}%
	\end{wrapfigure}%
	}]%
{Какую силу надо приложить к рычагу в точке $B$, чтобы он оставался в равновесии. Найти модуль силы и нарисовать вектор на рисунке.
}
%
{$F_B = \frac{8}{3}F$}

%%%%%%%%%%%%%%%%%%%%%%%%%%%%%%%%%%%%%%%%%%%%%%%%%%%%
\z[{\begin{wrapfigure}[2]{r}{0.2\textwidth}%
		\vspace{-15pt}%
		\begin{flushright}%
 			\includegraphics[width=0.2\textwidth]{5.1.2}%
		\end{flushright}%
	\end{wrapfigure}%
	}]%
{Масса первого груза $1\, \text{кг}$, масса третьего груза $2\, \text{кг}$. Какова масса второго груза?
}
%
{$m_2 = \frac{3m_3 - 4m_1}{2} = 1\ \text{кг}$}

%%%%%%%%%%%%%%%%%%%%%%%%%%%%%%%%%%%%%%%%%%%%%%%%%%%%
\z{Стержень массой $m = 9\, \text{кг}$ и длиной $l = 1\, \text{м}$ лежит на двух опорах. Одна из них находится у левого края стержня, а другая на расстоянии $a = 10\, \text{см}$ от правого края. С какой силой действует на стержень каждая из опор?
}
%
{$F_1 = \frac{mg(l - 2a)}{2(l - a)} = 40\ \text{Н}, \quad F_2 = \frac{mgl}{2(l - a)} = 50\, \text{Н}$}

%%%%%%%%%%%%%%%%%%%%%%%%%%%%%%%%%%%%%%%%%%%%%%%%%%%%
\z[{\begin{wrapfigure}[3]{r}{0.2\textwidth}%
		\vspace{-10pt}%
		\begin{flushright}%
 			\includegraphics[width=0.2\textwidth]{5.1.4}%
		\end{flushright}%
	\end{wrapfigure}%
	}]%
{Прямолинейный кусок проволоки массой $40\, \text{г}$ подвешен за середину (см. рисунок). Левую половину куска согнули, как показано на рисунке. Какой массы грузик надо подвесить в точке $A$, чтобы восстановить равновесие?
}
%
{$M = \frac{m}{4} = 10\, \text{г}$}

%%%%%%%%%%%%%%%%%%%%%%%%%%%%%%%%%%%%%%%%%%%%%%%%%%%%
\z[{\begin{wrapfigure}[7]{r}{0.3\textwidth}%
		\vspace{-10pt}%
		\begin{flushright}%
 			\includegraphics[width=0.3\textwidth]{5.1.5}%
		\end{flushright}%
	\end{wrapfigure}%
	}]%
{На двух вертикально расположенных пружинах, в недеформированном состоянии имеющих одинаковую длину, горизонтально подвешен невесомый стержень (см. рисунок). Жесткости пружин $k_1 = 0.02\, \text{Н/м}$ и $k_2 = 0.03\, \text{Н/м}$, расстояние между ними $d = 1\, \text{м}$. На каком расстоянии $x$ от первой пружины нужно подвесить груз, чтобы стержень остался в горизонтальном положении?
}
%
{$x = \frac{d k_2}{k_1 + k_2} = 0.6\, \text{м}$}

%%%%%%%%%%%%%%%%%%%%%%%%%%%%%%%%%%%%%%%%%%%%%%%%%%%%
\z[{\begin{wrapfigure}[3]{r}{0.3\textwidth}%
		\vspace{-15pt}%
		\begin{flushright}%
 			\includegraphics[width=0.3\textwidth]{5.1.6}%
		\end{flushright}%
	\end{wrapfigure}%
	}]%
{Линейка массы $0.01\, \text{кг}$ лежит на двух опорах так, как это показано на рисунке. На один конец линейки положен груз. Какова масса груза, при которой возможно равновесие?
}
%
{$m = 7.5\, \text{г}$
}

%%%%%%%%%%%%%%%%%%%%%%%%%%%%%%%%%%%%%%%%%%%%%%%%%%%%
\z[{\begin{wrapfigure}[3]{r}{0.2\textwidth}%
		\vspace{-10pt}%
		\begin{flushright}%
 			\includegraphics[width=0.2\textwidth]{5.1.7}%
		\end{flushright}%
	\end{wrapfigure}%
	}]%
{Стержень массой $m$ и длиной $l$ удерживают в горизонтальном положении с помощью двух точечных опор, расположенных на расстоянии $l/5$ друг от друга. Найти силы реакции опор, считая, что на одну из них стержень опирается самым краем.
}
%
{$N = \frac{5}{2}mg, \quad N_2 = \frac{3}{2}mg$}